
\medskip
\noindent
\SGARefTarget{sga7.sga7i.eix.n2-3}{2.3}. Suppose once again that $S$ is henselian and that
$\ell\ne p=\operatorname{char}k$.  We now use the inclusion
\begin{equation}\tag{\SGARefTarget{sga7.sga7i.eix.n2-3-1}{2.3.1}}
 T_\ell(T_o)\hookrightarrow T_\ell(A_o^o)
\end{equation}
of (\SGARefLink{sga7.sga7i.eix.n2-1-6}{2.1.6}), together with the equivalence of categories recalled in
\SGARefLink{sga7.sga7i.eix.n2-2-4}{2.2.4}.  We obtain a canonical twisted constant $\ell$-adic subsheaf
of $T_\ell(A^o)^f$, called the \underline{toric part of}
$T_\ell(A^o)$,
\begin{equation}\tag{\SGARefTarget{sga7.sga7i.eix.n2-3-2}{2.3.2}}
 T_\ell(A^o)^t\hookrightarrow T_\ell(A^o)^f
 \hookrightarrow T_\ell(A^o),
\end{equation}
characterized by the condition that its special fiber is the inclusion
(\SGARefLink{sga7.sga7i.eix.n2-3-1}{2.3.1}):
\begin{equation}\tag{\SGARefTarget{sga7.sga7i.eix.n2-3-3}{2.3.3}}
 T_\ell(A^o)^t\times_Ss=T_\ell(T_o)
 \hookrightarrow T_\ell(A_o^o).
\end{equation}
Later we shall give an interpretation of the toric part of
$T_\ell(A^o)$ as the Tate module $T_\ell(T)$ of a torus $T$ over $S$
lifting the torus $T_o$ over $k$.  That interpretation will allow the
definition of the toric part to be extended to the case $\ell=p$.
For the moment, and to fix ideas, we confine ourselves to the case
$\ell\ne p$.  In this case the toric part of $T_\ell(A^o)$ may be
described by the corresponding $\mathbb Z_\ell$-submodule of
$T_\ell(A(\overline K))$:
\begin{equation}\tag{\SGARefTarget{sga7.sga7i.eix.n2-3-4}{2.3.4}}
 T_\ell(A(\overline K))^t\hookrightarrow
 T_\ell(A(\overline K))^f\hookrightarrow
 T_\ell(A(\overline K)).
\end{equation}
It is stable under the action of $\pi$; its first term is defined as
$T_\ell(A^o)^t(\overline K)$ and is, as one would expect, again called
the \underline{toric submodule} of the Tate module
$T_\ell(A(\overline K))$.  Thus this Tate module carries the canonical
three-step filtration (\SGARefLink{sga7.sga7i.eix.n2-3-4}{2.3.4}).  Applying the same considerations to
the dual abelian scheme $A'_K$, we obtain the analogous three-step
filtration
\begin{equation}\tag{\SGARefTarget{sga7.sga7i.eix.n2-3-5}{2.3.5}}
 T_\ell(A'(\overline K))^t\hookrightarrow
 T_\ell(A'(\overline K))^f\hookrightarrow
 T_\ell(A'(\overline K)).
\end{equation}
The middle terms in (\SGARefLink{sga7.sga7i.eix.n2-3-4}{2.3.4}) and (\SGARefLink{sga7.sga7i.eix.n2-3-5}{2.3.5}), the ``fixed parts,'' are
described in terms of the $\pi$-module structures of the Tate modules
of $A(\overline K)$ and $A'(\overline K)$ as their submodules of
invariants under the inertia group $I$ (\SGARefLink{sga7.sga7i.eix.n2-2-5}{2.2.5}).  The toric parts are
then completely described in terms of the fixed parts and the pairing
$\varphi$ of (\SGARefLink{sga7.sga7i.eix.n1-0-2}{1.0.2}), by the following theorem (a special case of which
already appears in [11]).

\medskip
\noindent
\underline{Theorem} \SGARefTarget{sga7.sga7i.eix.n2-4}{2.4}. (\underline{Orthogonality theorem}.)
\underline{Let $S$ be a henselian trait, let $A_K$ be an abelian scheme
over its fraction field $K$, let $A'_K$ be the dual abelian scheme, and
let $\ell$ be a prime different from the residual characteristic of
$S$.  The Tate modules $T_\ell(A(\overline K))$ and
$T_\ell(A'(\overline K))$ then carry the canonical filtrations (\SGARefLink{sga7.sga7i.eix.n2-3-4}{2.3.4})
and (\SGARefLink{sga7.sga7i.eix.n2-3-5}{2.3.5}), respectively, by the ``fixed part'' and the ``toric
part.''  The toric part $T_\ell(A(\overline K))^t$ is the intersection
of the fixed part $T_\ell(A(\overline K))^f$ with the orthogonal of the
fixed part
$T_\ell(A'(\overline K))^f=T_\ell(A'(\overline K))^I$ with respect to
the canonical pairing $\varphi$ of (\SGARefLink{sga7.sga7i.eix.n1-0-3}{1.0.3}).  Symmetrically, the toric
part $T_\ell(A'(\overline K))^t$ is the intersection of
$T_\ell(A'(\overline K))^f$ with the orthogonal of the fixed part
$T_\ell(A(\overline K))^f=T_\ell(A(\overline K))^I$.  (Here the
superscript $I$ denotes the submodule of invariants under the action of
the inertia group $I$.)}

\smallskip
\noindent
\underline{Proof.}\quad a) Let us prove that
$T_\ell(A(\overline K))^t$ is orthogonal to
$T_\ell(A'(\overline K))^f$.  To do so, we interpret the restriction of
$\varphi$ to the fixed parts of the Tate modules using the canonical
biextension $W^o=W$ of (\SGARefLink{sga7.sga7i.eix.n1-4-4}{1.4.4}).  This is the canonical extension of the
biextension $W_K$ of $(A_K,A'_K)$ by $\mathbb G_{mK}$ to a biextension
of $(A^o,A'{}^o)$ by $\mathbb G_{mS}$.  By VIII~\SGARefLink{sga7.sga7i.eviii.n2-2}{2.2}, $W$ gives rise to
a pairing
\begin{equation}\tag{\SGARefTarget{sga7.sga7i.eix.n2-4-1}{2.4.1}}
 \varphi_W:T_\ell(A^o)\times T_\ell(A'{}^o)
 \longrightarrow T_\ell(\mathbb G_{mS}),
\end{equation}
from which the pairing $\varphi$ of (\SGARefLink{sga7.sga7i.eix.n1-0-3}{1.0.3}) is obtained by passing to
$\overline K$-valued points.  It follows that the restriction of
$\varphi$ to the fixed parts considered in \SGARefLink{sga7.sga7i.eix.n2-4}{2.4} is obtained, by passing
to $\overline K$-valued points, from the pairing
\begin{equation}\tag{\SGARefTarget{sga7.sga7i.eix.n2-4-2}{2.4.2}}
 T_\ell(A^o)^f\times T_\ell(A'{}^o)^f
 \longrightarrow T_\ell(\mathbb G_{mS}),
\end{equation}
the restriction of (\SGARefLink{sga7.sga7i.eix.n2-4-1}{2.4.1}).  Since the two factors on the left of
(\SGARefLink{sga7.sga7i.eix.n2-4-2}{2.4.2}) are now twisted constant $\ell$-adic sheaves, this pairing may
be studied on the special fibers.  The fibers of the two factors are
precisely the Tate modules $T_\ell(A_o^o)$ and
$T_\ell(A'{}_o^o)$ (\SGARefLink{sga7.sga7i.eix.n2-2-3-3}{2.2.3.3}), and the pairing induced on special fibers
by (\SGARefLink{sga7.sga7i.eix.n2-4-2}{2.4.2}) is
\begin{equation}\tag{\SGARefTarget{sga7.sga7i.eix.n2-4-3}{2.4.3}}
 T_\ell(A_o^o)\times T_\ell(A'{}_o^o)
 \longrightarrow T_\ell(\mathbb G_{mk}),
\end{equation}
obtained from the biextension $W_o$ of $(A_o^o,A'{}_o^o)$ by
$\mathbb G_{mk}$, the restriction of $W$ to the special fibers.

Recall from \SGARefLink{sga7.sga7i.eix.n2-3}{2.3} that the toric part of
$T_\ell(A(\overline K))$ is obtained by passing to
$\overline K$-valued points in the twisted constant $\ell$-adic
subsheaf $T_\ell(A^o)^t$ of $T_\ell(A^o)^f$.  Thus the orthogonality
relation to be proved says that the pairing induced by (\SGARefLink{sga7.sga7i.eix.n2-4-2}{2.4.2}),
\begin{equation}\tag{\SGARefTarget{sga7.sga7i.eix.n2-4-4}{2.4.4}}
 T_\ell(A^o)^t\times T_\ell(A'{}^o)^f
 \longrightarrow T_\ell(\mathbb G_{mS}),
\end{equation}
is identically zero.  It is enough to verify this for the corresponding
pairing on the special fibers, which by the definition of the toric
part is
\begin{equation}\tag{\SGARefTarget{sga7.sga7i.eix.n2-4-5}{2.4.5}}
 T_\ell(T_o)\times T_\ell(A'{}_o^o)
 \longrightarrow T_\ell(\mathbb G_{mk}),
\end{equation}
induced by (\SGARefLink{sga7.sga7i.eix.n2-4-3}{2.4.3}).  By VIII~\SGARefLink{sga7.sga7i.eviii.n4-10}{4.10} this pairing is necessarily zero,
which proves the asserted orthogonality.

\smallskip
\noindent
b) Let us prove that the inclusion
\begin{equation}\tag{\SGARefTarget{sga7.sga7i.eix.n2-4-6}{2.4.6}}
 T_\ell(A(\overline K))^t\subset
 T_\ell(A(\overline K))^f\cap
 \bigl(T_\ell(A'(\overline K))^f\bigr)^\perp
\end{equation}
is an equality, where $\perp$ denotes the orthogonal.  Equivalently, in
the pairing (\SGARefLink{sga7.sga7i.eix.n2-4-2}{2.4.2}), the annihilator of the second factor---which, as
we have just seen, contains the toric part $T_\ell(A^o)^t$---is equal
to that toric part.  Since these are again relations between twisted
constant $\ell$-adic sheaves, it is enough to check the assertion on
the special fibers: namely, that the inclusion
\begin{equation}\tag{\SGARefTarget{sga7.sga7i.eix.n2-4-7}{2.4.7}}
 T_\ell(T_o)\subset T_\ell(A'{}_o^o)^\perp
\end{equation}
is an equality.  In other words, we must verify that the biextension
$W_o$ of $(A_o^o,A'{}_o^o)$ is \underline{nondegenerate on the left}
(VIII~\SGARefLink{sga7.sga7i.eviii.n4-11}{4.11}).  But this is precisely what we already observed in \SGARefLink{sga7.sga7i.eix.n1-5}{1.5},
as a consequence of a result from Raynaud's thesis.
